\documentclass{article}
\usepackage{fullpage}
\begin{document}
\noindent\textbf{Describe your professional, academic, or life experience and skills you have that enhance your ability to be an excellent contributing member to the SULI Program.}
\\\\
My technical depth, research experience in computational science, 
and interdisciplinary communication skills equip me to be a strong contributor 
to the SULI program. Throughout my academic journey, I have learned how to work independently 
in research environments, collaborate with people from diverse disciplines, and 
communicate complex scientific ideas clearly, which are paramount for 
being productive in a national laboratory setting.
\\

Through research experience, coursework, and personal projects, I have become fluent in 
Python, C, C++, LaTeX, and I am comfortable using scientific computing packages such as 
NumPy, SciPy, Qiskit, Cirq, and Matplotlib. I have also worked in high-performance 
computing environments and gained familiarity with software development 
practices through open-source contributions, including 
version control, CI/CD workflows, and modular code design. These experiences have 
taught me how to write efficient, reproducible scientific code, which is essential for implementing 
computational methods across various scientific domains.
\\

Beyond programming, I have learned how to approach and 
communicate complex scientific problems. I first became interested in computational 
methods through Monte Carlo algorithms during my first semester of college. I was intrigued
by the versatility of these methods, and used them extensively in enhanced molecular 
dynamics sampling and optimizing quantum circuit simulator protocols.  Learning LaTeX
and practicing technical writing has strengthened my ability to communicate scientific 
ideas, which has been useful for preparing materials for conferences such as ACS Fall 2025 and for 
discussing results with experimental collaborators. Through these experiences, I have learned how to 
adapt my communication to audiences with different technical backgrounds, which is crucial for interdisciplinary 
research. 
\\

Working with computationally intensive systems such as large biomolecular systems or 
quantum circuits has taught me how to evaluate tradeoffs between numerical methods and 
choose appropriate approaches for different problems. Collaborating with 
postdocs and my principal investigators, and exploring relevant literature 
has helped me learn how to incorporate feedback and make decisions about the 
feasibility of different computational strategies. 
\\

My experience in quantum algorithms, molecular dynamics simulations, 
and computational physics has given me a broad perspective on both the utility and 
limitations of current computational methods. Conducting research in different fields has 
taught me about the importance of approaching problems from multiple angles and connecting 
methodologies across disciplines. I am eager to contribute my technical versatility and interdisciplinary knowledge 
to computational projects through the SULI program.


\end{document}